\chapter{Анализ спектров дейтерированных изотопологов сероводорода в области полосы $\nu_2$}\label{ch:ch3}

В предыдущей главе обратная спектроскопическая задача рассматривалась на уровне колебательных центров
полос молекул высокой симметрии. Настоящая глава продолжает это исследование применительно к
высокоразрешенным колебательно-вращательным спектрам молекул типа асимметричного волчка. Основное
внимание уделено фундаментальной деформационной полосе $\nu_2$ дейтерированных изотопологов
сероводорода D$_2{}^{32}$S, HD$^{32}$S и HD$^{34}$S.

В отличие от анализа центров полос, рассматриваемая задача включает последовательное отнесение
индивидуальных переходов, восстановление энергетической структуры возбужденного состояния $(010)$,
определение параметров эффективного гамильтониана и анализ абсолютных интенсивностей линий. Особое
значение имеет учет водород-дейтериевого обмена в экспериментальном образце, вследствие которого
парциальные давления отдельных изотопологов не могли быть заданы непосредственно и определялись с
использованием соотношений теории изотопозамещения.

Материал главы объединяет два взаимосвязанных этапа исследования, выполненных по одним и тем же
экспериментальным спектрам. Сначала рассматривается полоса $\nu_2$ молекулы D$_2{}^{32}$S, после чего
аналогичный подход применяется к монодейтерированным изотопологам HD$^{32}$S и HD$^{34}$S. Такое
изложение позволяет последовательно сопоставить особенности спектров молекул симметрии C$_{2v}$ и
C$_s$, а также оценить применимость единой методики к определению положений и интенсивностей линий.

\section{Анализ спектра молекулы D$_2{}^{32}$S в области полосы $\nu_2$}\label{sec:ch3/D2S}\label{sec:ch2/D2S}

В данном параграфе рассматривается молекула дидейтерированного сероводорода D$_2{}^{32}$S, для
которой была решена обратная спектроскопическая задача в области фундаментальной деформационной
полосы $\nu_2$. Под обратной задачей здесь понимается восстановление параметров эффективного
гамильтониана по экспериментально проинтерпретированным положениям линий, уточнение энергетической
структуры состояния $(010)$, а также определение параметров эффективного оператора дипольного момента
по измеренным абсолютным интенсивностям линий. Основные результаты этого анализа опубликованы в
работе~\cite{d2s_sydow2019}.

Актуальность исследования D$_2$S связана с тем, что сероводород и его изотопологи представляют интерес
для атмосферной оптики, планетологии, астрофизики и задач, связанных с серным циклом
\cite{h2s_farquhar2000_science,h2s_farquhar2000_nature,h2s_thaddeus1972,h2s_vastel2003}. Кроме того,
молекула H$_2$S относится к числу легких молекул типа асимметричного волчка, для которых
колебательно-вращательные взаимодействия проявляются особенно заметно. Поэтому данные по
дейтерированным изотопологам важны не только как самостоятельная спектроскопическая информация, но и
как дополнительное ограничение при уточнении поверхности потенциальной энергии и поверхности
дипольного момента. При этом для D$_2$S до настоящего анализа существовало ограниченное число работ по
положениям линий в ИК-диапазоне \cite{d2s_gillis1985,d2s_camy1985,d2s_camy1988}, а данные об
абсолютных интенсивностях были опубликованы только для растягивающих полос $\nu_1$ и $\nu_3$
\cite{d2s_lamouroux2008}. Поэтому анализ фундаментальной полосы $\nu_2$ был необходим как для
расширения энергетической информации, так и для построения согласованного набора интенсивностей.

\subsection{Экспериментальные условия и общий вид спектра}

Спектры D$_2$S были зарегистрированы на Фурье-спектрометре Bruker IFS~125HR в лаборатории
Технического университета Брауншвейга. Исходная область регистрации составляла 600--3000 см$^{-1}$,
однако для настоящего анализа использовался диапазон 750--1200 см$^{-1}$, в котором локализована
полоса $\nu_2$. Эксперимент выполнялся для коммерческого образца D$_2$S с заявленной чистотой 97~\%.
Основные условия регистрации приведены в табл.~\ref{tab:d2s_experiment}.

\begin{table}[htbp]
\centering
\caption{Условия регистрации спектров D$_2{}^{32}$S, использованных при анализе полосы $\nu_2$.}
\label{tab:d2s_experiment}
\begingroup
\scriptsize
\setlength{\tabcolsep}{3pt}
\renewcommand{\arraystretch}{1.15}
\begin{tabular}{|c|c|c|c|c|c|c|c|c|c|}
\hline
Спектр &
\shortstack{Разрешение,\\см$^{-1}$} &
\shortstack{Число\\сканов} &
\shortstack{Диапазон,\\см$^{-1}$} &
Детектор &
Делитель &
\shortstack{$L$,\\м} &
\shortstack{Апертура,\\мм} &
\shortstack{$T$,\\$^\circ$C} &
\shortstack{$P$,\\Па} \\
\hline
I  & 0.003 & 1200 & 600--2200 & MCT & KBr & 163 & 1.70 & $23\pm0.8$ & 450 \\
II & 0.002 & 960  & 600--3000 & MCT & KBr & 4   & 1.15 & $23\pm0.8$ & 450 \\
\hline
\end{tabular}
\endgroup
\end{table}

На рис.~\ref{fig:d2s_overview} показан общий вид экспериментальных спектров и рассчитанный спектр в
области полосы $\nu_2$. Видно, что в этой области хорошо выражены $P$- и $R$-ветви. Более детальные
фрагменты спектра, использованные при отнесении линий, приведены на рис.~\ref{fig:d2s_fragments}. На
фрагменте $P$-ветви неотнесенные линии в основном связаны с изотопологом D$_2{}^{34}$S, тогда как в
области $Q/R$-ветвей заметен вклад линий HD$^{32}$S. Это обстоятельство связано с быстрым
водород-дейтериевым обменом в образце: в течение длительной регистрации часть молекул D$_2$S
переходила в HDS и H$_2$S из-за присутствия остаточных следов воды на стенках кюветы и окнах.

\begin{figure}[htbp]
    \centering
    \maybeincludegraphics[width=0.90\textwidth]{images/d2s_fig1.pdf}
    \caption{Общий вид экспериментальных спектров D$_2{}^{32}$S в области 750--1200 см$^{-1}$ и
    соответствующий рассчитанный спектр.}
    \label{fig:d2s_overview}
\end{figure}

\begin{figure}[htbp]
    \centering
    \maybeincludegraphics[width=0.49\textwidth]{images/d2s_fig2.pdf}\hfill
    \maybeincludegraphics[width=0.49\textwidth]{images/d2s_fig3.pdf}
    \caption{Фрагменты высокоразрешенного спектра D$_2{}^{32}$S в области $P$-ветви (слева) и
    $Q/R$-ветвей (справа) полосы $\nu_2$.}
    \label{fig:d2s_fragments}
\end{figure}

Интенсивность отдельной линии определялась из закона Бугера--Ламберта--Бера через площадь линии
поглощения $A_{\text{line}}$, парциальное давление исследуемого газа $P$, температуру $T$ и оптическую
длину пути $L$:
\begin{equation}
S=\frac{1}{\lg e}\frac{k_B T}{PL}A_{\text{line}},
\label{eq:d2s_beer_lambert}
\end{equation}
где
\begin{equation}
A_{\text{line}}=\int \lg\left(\frac{I_0}{I}\right)_{\widetilde{\nu}}d\widetilde{\nu}.
\label{eq:d2s_line_area}
\end{equation}
Профили отдельных линий описывались с использованием профиля Хартмана--Трана, рекомендованного для
высокоточного описания изолированных линий высокого разрешения \cite{d2s_tennyson2014}. Для анализа
интенсивностей использовался спектр, записанный с короткой оптической длиной пути 4~м, поскольку в
этом случае уменьшается риск насыщения сильных линий и повышается надежность определения площади
линии.

\subsection{Отнесение линий и энергетическая структура состояния $(010)$}

Молекула D$_2$S является асимметричным волчком с параметром асимметрии
$\kappa=(2B-A-C)/(A-C)\simeq0.357$. Нормальные координаты $q_1$ и $q_2$ имеют симметричный тип, а
координата $q_3$ --- антисимметричный тип. Поэтому состояние $(010)$ является симметричным, а полоса
$\nu_2$ относится к полосам $b$-типа. Для таких полос выполняются правила отбора
\begin{equation}
\Delta J=0,\pm1,\qquad \Delta K_a=\text{нечетное},\qquad \Delta K_c=\text{нечетное}.
\label{eq:d2s_selection_rules}
\end{equation}

Отнесение переходов выполнялось методом комбинационных разностей в основном колебательном состоянии.
В качестве исходного приближения использовались опубликованные вращательные параметры основного
состояния \cite{d2s_camy1985}. В ходе анализа было установлено, что при увеличении квантовых чисел
$J$ и $K_a$ различия между экспериментальными комбинационными разностями и значениями, рассчитанными
по литературным параметрам, систематически возрастают. Поэтому параметры основного состояния были
уточнены: к исходному набору микроволновых и ИК-данных были добавлены 270 экспериментальных
комбинационных разностей, полученных в настоящей работе. Новый набор параметров воспроизводит эти
комбинационные разности со среднеквадратичным отклонением $d_{\textrm{rms}}=1.4\times10^{-4}$
см$^{-1}$, что примерно в 2.9 раза лучше, чем воспроизведение тем же набором данных с использованием
предыдущих параметров. Уточненные параметры приведены в табл.~\ref{tab:d2s_ground_parameters}.

\begin{table}[htbp]
\centering
\caption{Уточненные вращательные параметры основного колебательного состояния D$_2$S, см$^{-1}$.}
\label{tab:d2s_ground_parameters}
\small
\begin{tabular}{|c|c|c|}
\hline
Параметр & \cite{d2s_camy1985} & Настоящий анализ \\
\hline
$A$ & 5.49167887 & 5.49167798(24) \\
$B$ & 4.512610865 & 4.51260976(17) \\
$C$ & 2.444178310 & 2.44417753(12) \\
$\Delta_K\times10^3$ & 0.9699746 & 0.970034(16) \\
$\Delta_{JK}\times10^3$ & -0.5623655 & -0.562393(11) \\
$\Delta_J\times10^3$ & 0.1601221 & 0.1600954(35) \\
$\delta_K\times10^3$ & -0.0226215 & -0.0226046(61) \\
$\delta_J\times10^3$ & 0.07195620 & 0.0719499(17) \\
$H_K\times10^6$ & 0.20828 & 0.21089(49) \\
$H_{KJ}\times10^6$ & 0.13616 & 0.13494(40) \\
$H_{JK}\times10^6$ & -0.181561 & -0.181317(93) \\
$H_J\times10^6$ & 0.032592 & 0.032443(26) \\
$d_{\textrm{rms}}\times10^4$ & 4.1 & 1.4 \\
\hline
\end{tabular}
\end{table}


После уточнения основного состояния к полосе $\nu_2$ было отнесено 1742 перехода с максимальными
значениями квантовых чисел $J_{\max}=30$ и $K_{a,\max}=21$. Из этих переходов было получено 535
колебательно-вращательных энергий состояния $(010)$, что примерно в два раза превышает объем данных,
использованный в предыдущих исследованиях. Сводная статистика отнесения приведена в
табл.~\ref{tab:d2s_statistics}.

\begin{table}[htbp]
\centering
\caption{Статистика анализа полосы $\nu_2$ молекулы D$_2{}^{32}$S.}
\label{tab:d2s_statistics}
\small
\begin{tabular}{|c|c|c|c|c|c|c|c|c|}
\hline
Источник & Центр, см$^{-1}$ & $J_{\max}$ & $K_{a,\max}$ & $N_{tr}$ & $N_l$ & $m_1$, \% & $m_2$, \% & $m_3$, \% / $m_4$, \% \\
\hline
\cite{d2s_gillis1985} & 855.40475 & 21 & 17 & 655 & -- & -- & -- & -- \\
\cite{d2s_camy1988} & 855.404206 & 23 & 16 & -- & 284 & -- & -- & -- \\
Настоящий анализ & 855.404105 & 30 & 21 & 1742 & 535 & 76.1 & 13.3 & 6.2 / 4.4 \\
\hline
\end{tabular}
\end{table}


Для описания состояния $(010)$ использовался эффективный гамильтониан Ватсона в $A$-редукции и
$I^r$-представлении \cite{0102}. Поскольку состояние $(010)$ в рассматриваемой области можно считать
изолированным, подгонка выполнялась без явного учета резонансных взаимодействий с другими
колебательными состояниями. Взвешенная подгонка 535 экспериментальных энергетических значений дала
набор из 34 спектроскопических параметров. Полученный набор воспроизводит исходные энергии со
среднеквадратичным отклонением $d_{\textrm{rms}}=1.31\times10^{-4}$ см$^{-1}$. Основные параметры
состояния $(010)$ и сравнение с предыдущими работами приведены в
табл.~\ref{tab:d2s_010_parameters}, а распределение невязок показано на рис.~\ref{fig:d2s_residuals}.

\begin{table}[htbp]
\centering
\caption{Основные спектроскопические параметры состояния $(010)$ D$_2$S, см$^{-1}$.}
\label{tab:d2s_010_parameters}
\small
\begin{tabular}{|c|c|c|c|}
\hline
Параметр & \cite{d2s_gillis1985} & \cite{d2s_camy1988} & Настоящий анализ \\
\hline
$E$ & 855.40475 & 855.404206 & 855.404105(25) \\
$A$ & 5.6248675 & 5.62502454 & 5.6250158(28) \\
$B$ & 4.58662638 & 4.5863907 & 4.5866339(17) \\
$C$ & 2.42144557 & 2.4215077 & 2.42144601(73) \\
$\Delta_K\times10^3$ & 1.120229 & 1.12682429 & 1.125781(75) \\
$\Delta_{JK}\times10^3$ & -0.639135 & -0.64172117 & -0.641256(48) \\
$\Delta_J\times10^3$ & 0.17468 & 0.177659290 & 0.177675(11) \\
$\delta_K\times10^3$ & -0.009473 & -0.00044164 & -0.000217(21) \\
$\delta_J\times10^3$ & 0.078936 & 0.080781116 & 0.0807954(53) \\
$H_K\times10^6$ & 0.3242 & 0.3274980 & 0.32902(69) \\
$H_{KJ}\times10^6$ & 0.0450 & 0.1116061 & 0.09615(69) \\
$H_{JK}\times10^6$ & -0.18044 & -0.2119880 & -0.20371(24) \\
$H_J\times10^6$ & 0.012466 & 0.04085927 & 0.040684(39) \\
\hline
\end{tabular}
\end{table}


\begin{figure}[htbp]
    \centering
    \maybeincludegraphics[width=\textwidth,trim=0 6.5cm 0 0,clip]{images/d2s_fig4.pdf}
    \caption{Остаточные отклонения <<эксперимент--расчет>> для положений линий, энергий состояния
    $(010)$ и интенсивностей линий полосы $\nu_2$ молекулы D$_2{}^{32}$S.}
    \label{fig:d2s_residuals}
\end{figure}

\subsection{Оценка парциального давления D$_2$S в образце}

Особенность данного эксперимента состояла в том, что общее давление, температура и длина оптического
пути контролировались с высокой точностью, однако парциальное давление D$_2$S в течение записи
оставалось неизвестным из-за водород-дейтериевого обмена. Наличие в образце HDS и H$_2$S подтверждается
фрагментами спектра, приведенными на рис.~\ref{fig:d2s_hds_h2s}. Поэтому для корректного определения
абсолютных интенсивностей потребовалась дополнительная процедура оценки доли D$_2$S в газовой смеси.

\begin{figure}[htbp]
    \centering
    \maybeincludegraphics[width=0.82\textwidth]{images/d2s_fig5.pdf}
    \caption{Фрагменты спектра, подтверждающие наличие в образце линий H$_2$S и HDS, возникших в
    результате водород-дейтериевого обмена.}
    \label{fig:d2s_hds_h2s}
\end{figure}

Процедура оценки была основана на теории изотопозамещения \cite{byk,Ulenikov2019a,Ulenikov2020a}. Так
как дипольный момент молекулы как системы ядер и электронов не зависит от масс ядер, параметры
эффективного дипольного момента изотополога можно связать с параметрами материнской молекулы. Для
главного параметра полосы $\nu_2$ молекулы D$_2$S получено соотношение
\begin{equation}
{}^{010}\widetilde{\mu}_{x1}= {}^{100}\mu_{x1}\alpha^2_1
\left(\frac{\omega_1}{\widetilde{\omega}_2}\right)^{1/2}+
{}^{010}\mu_{x1}\alpha^2_2
\left(\frac{\omega_2}{\widetilde{\omega}_2}\right)^{1/2}.
\label{eq:d2s_mu_iso}
\end{equation}
При использовании параметров H$_2$S \cite{h2s_ulenikov2018_bending,h2s_lechuga1984} это соотношение
дает значение ${}^{010}\widetilde{\mu}_{x1}=0.008519$~D. Оно было использовано для расчета
интенсивностей 16 изолированных пробных линий с малыми значениями $J$ и $K_a$. Парциальное давление
оценивалось по формуле
\begin{equation}
P_{\text{part}}=\frac{{}^{(\text{prob})}S^N_{\nu}}{{}^{(\text{calc})}S^N_{\nu}}P_{\text{samp}},
\label{eq:d2s_partial_pressure}
\end{equation}
где $P_{\text{samp}}=4.5$ мбар --- полное давление образца. Первичная оценка дала
$P_{\text{part}}=2.40\pm0.27$ мбар, однако остатки для линий $P$- и $R$-ветвей показали необходимость
учесть следующий по значимости параметр эффективного дипольного момента ${}^{010}\widetilde{\mu}_{x4}$.

На следующем шаге для шести пробных значений давления были заново определены интенсивности 16 линий и
выполнена подгонка двух параметров ${}^{010}\widetilde{\mu}_{x1}$ и
${}^{010}\widetilde{\mu}_{x4}$. Пересечение рассчитанной зависимости с теоретически предсказанным
значением ${}^{010}\widetilde{\mu}_{x1}$ дало оптимальное отношение
${}^{(\text{calc})}S^N_{\nu}/{}^{(\text{prob})}S^N_{\nu}=1.892\pm0.028$ и, соответственно,
$P_{\text{part}}=2.378\pm0.036$ мбар. Результаты этой процедуры приведены в
табл.~\ref{tab:d2s_pressure_two_parameters} и на рис.~\ref{fig:d2s_pressure}.

\begin{table}[htbp]
\centering
\caption{Оценка парциального давления D$_2$S по 16 пробным переходам с учетом параметров ${}^{010}\widetilde{\mu}_{x1}$ и ${}^{010}\widetilde{\mu}_{x4}$; $R_S=\langle{}^{(\textrm{calc})}S^N_\nu/{}^{(\textrm{prob})}S^N_\nu\rangle$.}
\label{tab:d2s_pressure_two_parameters}
\begingroup
\scriptsize
\setlength{\tabcolsep}{3pt}
\renewcommand{\arraystretch}{1.15}
\begin{tabular}{|p{0.22\textwidth}|c|c|c|p{0.25\textwidth}|}
\hline
Группа переходов & Число линий & $R_S$ & $P_{\textrm{part}}$, мбар & Примечание \\
\hline
Все пробные линии & 16 & $1.892\pm0.028$ & $2.378\pm0.036$ & оптимальная оценка \\
$P$-ветвь & 7 & 2.089 & 2.154 & оценка без уточнения ${}^{010}\widetilde{\mu}_{x4}$ \\
$R$-ветвь & 9 & 1.714 & 2.625 & оценка без уточнения ${}^{010}\widetilde{\mu}_{x4}$ \\
\hline
\end{tabular}
\endgroup
\end{table}

\begin{figure}[htbp]
    \centering
    \maybeincludegraphics[width=0.82\textwidth]{images/d2s_fig6.pdf}
    \caption{Определение парциального давления D$_2$S по зависимости параметра
    ${}^{010}\widetilde{\mu}_{x1}$ от отношения рассчитанных и экспериментально определенных
    интенсивностей пробных линий.}
    \label{fig:d2s_pressure}
\end{figure}

\subsection{Анализ интенсивностей и параметры эффективного дипольного момента}

После определения парциального давления D$_2$S были получены экспериментальные интенсивности 280 линий
полосы $\nu_2$, соответствующие 324 переходам с $J_{\max}=16$ и $K_{a,\max}=9$. Площади линий
извлекались из подгонки профилей Хартмана--Трана; примеры такой подгонки приведены на
рис.~\ref{fig:d2s_lineshape}. Нормированная оптическая толщина при этом записывалась как
\begin{equation}
\tau(\nu)^{\textrm{exp}}=\frac{1}{\lg e}\lg\left(\frac{I_0(\nu)}{I(\nu)}\right).
\label{eq:d2s_tau}
\end{equation}

\begin{figure}[htbp]
    \centering
    \maybeincludegraphics[width=0.90\textwidth]{images/d2s_fig7.pdf}
    \caption{Примеры подгонки профилей отдельных линий полосы $\nu_2$ молекулы D$_2{}^{32}$S с
    использованием профиля Хартмана--Трана.}
    \label{fig:d2s_lineshape}
\end{figure}

Экспериментальные интенсивности были использованы в подгонке параметров эффективного оператора
дипольного момента. В итоговый набор вошло восемь параметров, приведенных в
табл.~\ref{tab:d2s_dipole_parameters}. Этот набор воспроизводит исходные 280 интенсивностей со
среднеквадратичным отклонением $d_{\textrm{rms}}=2.5$~\%. С учетом экспериментальной неопределенности
интенсивностей итоговая погрешность рассчитанных значений оценивается примерно как 3.5~\%.

\begin{table}[htbp]
\centering
\caption{Параметры эффективного дипольного момента полосы $\nu_2$ D$_2{}^{32}$S, D.}
\label{tab:d2s_dipole_parameters}
\small
\begin{tabular}{|c|c|c|}
\hline
Оператор & Параметр & Значение \\
\hline
$k_{zx}$ & ${}^{010}\widetilde{\mu}_{x1}\times10^{-2}$ & 0.8549(12) \\
$\{k_{zx},J^2\}$ & ${}^{010}\widetilde{\mu}_{x2}$ & -- \\
$\{k_{zx},J_z^2\}$ & ${}^{010}\widetilde{\mu}_{x3}\times10^{-5}$ & -0.821(56) \\
$\{ik_{zy},J_z\}$ & ${}^{010}\widetilde{\mu}_{x4}\times10^{-3}$ & 0.6367(38) \\
$\{k_{zz},iJ_y\}$ & ${}^{010}\widetilde{\mu}_{x5}\times10^{-4}$ & 0.1328(67) \\
$\{k_{zz},\{J_x,J_z\}\}$ & ${}^{010}\widetilde{\mu}_{x6}\times10^{-5}$ & 0.212(30) \\
$\frac{1}{2}[\{k_{zx},J_{xy}^2\}-\{ik_{zy},i\{J_x,J_y\}\}]$ & ${}^{010}\widetilde{\mu}_{x7}\times10^{-5}$ & -0.336(27) \\
$\frac{1}{2}[\{k_{zx},J_{xy}^2\}+\{ik_{zy},i\{J_x,J_y\}\}]$ & ${}^{010}\widetilde{\mu}_{x8}\times10^{-5}$ & 0.422(24) \\
\hline
\end{tabular}
\end{table}


Итогом анализа стал список из 1742 переходов полосы $\nu_2$ D$_2{}^{32}$S, для которых были получены
положения линий, энергетические значения верхнего состояния и рассчитанные индивидуальные
интенсивности. Таким образом, для фундаментальной деформационной полосы D$_2$S был построен
согласованный набор параметров эффективного гамильтониана и эффективного дипольного момента,
позволяющий воспроизводить как положения линий, так и их абсолютные интенсивности с точностью,
сопоставимой с экспериментальной.

\section{Анализ спектров HD$^{32}$S и HD$^{34}$S в области полосы $\nu_2$}\label{sec:ch3/HDS}\label{sec:ch2/HDS}

В данном параграфе рассматриваются монодейтерированные изотопологи сероводорода HD$^{32}$S и
HD$^{34}$S, для которых была решена обратная спектроскопическая задача в области фундаментальной
деформационной полосы $\nu_2$. Под обратной задачей здесь понимается отнесение экспериментальных
переходов, определение колебательно-вращательных энергий состояния $(010)$, уточнение параметров
эффективного гамильтониана Ватсона, а также определение параметров эффективного оператора дипольного
момента по измеренным интенсивностям линий. Основные результаты анализа опубликованы в
работе~\cite{hds_sydow2019}.

Этот анализ является естественным продолжением исследования полосы $\nu_2$ молекулы D$_2{}^{32}$S,
описанного в предыдущем параграфе. В обоих случаях использовались одни и те же экспериментальные
спектры, однако в настоящем параграфе объектом анализа являются линии HDS, возникшие в образце в
результате водород-дейтериевого обмена. Информация о спектрах HDS важна по тем же причинам, что и
данные по другим изотопологам сероводорода: сероводород участвует в серном цикле, встречается в
атмосферах планет и межзвездной среде, а его изотопологи дают дополнительные ограничения при уточнении
поверхности потенциальной энергии и поверхности дипольного момента
\cite{h2s_farquhar2000_science,h2s_farquhar2000_nature,h2s_thaddeus1972,h2s_vastel2003}. Ранее полоса
$\nu_2$ молекулы HD$^{32}$S уже анализировалась в работах~\cite{hds_camy1989,hds_ulenikov1995}, однако
число отнесенных переходов и диапазон квантовых чисел были существенно меньше, а систематический
анализ абсолютных интенсивностей для этой области отсутствовал.

\subsection{Экспериментальные условия и общий вид спектра}

Спектры были зарегистрированы на Фурье-спектрометре Bruker IFS~125HR в лаборатории Технического
университета Брауншвейга. Исходная область регистрации составляла 600--3000 см$^{-1}$, а для анализа
HDS использовался диапазон 690--1510 см$^{-1}$, включающий фундаментальную полосу $\nu_2$. Основные
экспериментальные условия приведены в табл.~\ref{tab:hds_experiment}. В отличие от обычного
эксперимента с заданным составом газа, в данном случае HDS образовывался непосредственно в кювете:
часть молекул D$_2$S переходила в HDS и H$_2$S из-за остаточных следов воды на стенках кювет и окнах.
Поэтому для количественного анализа интенсивностей потребовалось отдельно оценить парциальное давление
HD$^{32}$S в образце.

\begin{table}[htbp]
\centering
\caption{Условия регистрации спектров HDS в анализируемой области 750--1400 см$^{-1}$.}
\label{tab:hds_experiment}
\small
\resizebox{\textwidth}{!}{%
\begin{tabular}{c c c c c c c c c c c}
\hline
Спектр & Разрешение, см$^{-1}$ & Число сканов & Область, см$^{-1}$ & Детектор & Делитель &
$L$, м & Диафрагма, мм & $T$, $^\circ$C & $P$, Па & Калибровка \\
\hline
I  & 0.003 & 1200 & 600--2200 & MCT & KBr & 163 & 1.70 & $23\pm0.8$ & 450 & CO$_2$, N$_2$O \\
II & 0.002 &  960 & 600--3000 & MCT & KBr &   4 & 1.15 & $23\pm0.8$ & 450 & CO$_2$, N$_2$O \\
\hline
\end{tabular}%
}
\end{table}


На рис.~\ref{fig:hds_overview} показаны обзорные экспериментальные спектры и рассчитанный спектр HDS в
области полосы $\nu_2$. В спектре хорошо выражены $P$- и $R$-ветви, а в левой части видны линии
D$_2$S, относящиеся к тому же исходному образцу. Фрагменты спектра, использованные при отнесении
линий HD$^{32}$S, представлены на рис.~\ref{fig:hds_fragments}. Неотнесенные линии на этих фрагментах в
основном принадлежат HD$^{34}$S или горячим полосам.

\begin{figure}[htbp]
    \centering
    \maybeincludegraphics[width=0.90\textwidth]{images/hds_fig1.pdf}
    \caption{Общий вид экспериментальных спектров HDS в области полосы $\nu_2$ и соответствующий
    рассчитанный спектр.}
    \label{fig:hds_overview}
\end{figure}

\begin{figure}[htbp]
    \centering
    \maybeincludegraphics[width=0.49\textwidth]{images/hds_fig2.pdf}\hfill
    \maybeincludegraphics[width=0.49\textwidth]{images/hds_fig3.pdf}
    \caption{Фрагменты высокоразрешенного спектра HD$^{32}$S в областях $P$-ветви (слева) и
    $R$-ветви (справа) полосы $\nu_2$.}
    \label{fig:hds_fragments}
\end{figure}

\subsection{Отнесение линий и параметры состояния $(010)$ молекулы HD$^{32}$S}

Молекула HDS является асимметричным волчком с симметрией, изоморфной точечной группе C$_s$, и
параметром асимметрии $\kappa=(2B-A-C)/(A-C)\simeq -0.477$. Все три нормальные координаты $q_1$,
$q_2$ и $q_3$ имеют симметричный тип, поэтому полосы HDS являются гибридными. Для таких полос
использовались правила отбора
\begin{equation}
\Delta J=0,\pm1,\qquad \Delta K_a=\text{любое},\qquad \Delta K_c=\text{нечетное}.
\label{eq:hds_selection_rules}
\end{equation}

Отнесение переходов выполнялось методом комбинационных разностей в основном колебательном состоянии.
Энергии основного состояния были взяты из работы~\cite{d2s_camy1985}, где были определены
вращательные постоянные для D$_2{}^{32}$S, D$_2{}^{34}$S, HD$^{32}$S и HD$^{34}$S. В результате к полосе
$\nu_2$ молекулы HD$^{32}$S было отнесено 2684 перехода с максимальными квантовыми числами
$J_{\max}=25$ и $K_{a,\max}=17$. Для сравнения, в работе~\cite{hds_ulenikov1995} было отнесено около
1400 переходов с $J_{\max}=22$ и $K_{a,\max}=10$. Сводная статистика анализа для HD$^{32}$S и
HD$^{34}$S приведена в табл.~\ref{tab:hds_statistics}.

\begin{table}[htbp]
\centering
\caption{Статистическая информация для полосы $\nu_2$ молекул HD$^{32}$S и HD$^{34}$S.}
\label{tab:hds_statistics}
\begingroup
\scriptsize
\setlength{\tabcolsep}{2pt}
\renewcommand{\arraystretch}{1.05}
\begin{tabular}{@{}lccccccccc@{}}
\hline
Изотополог и источник &
Центр, см$^{-1}$ &
$J_{\max}$ &
$K_{a,\max}$ &
$N_{tr}$ &
$N_l$ &
$m_1$, \% &
$m_2$, \% &
$m_3$, \% &
$m_4$, \% \\
\hline
HD$^{32}$S, \cite{hds_ulenikov1995} & 1032.715171 & 22 & 10 & 1400 & 226 & --   & --   & --  & --  \\
HD$^{32}$S, наст. работа             & 1032.715193 & 25 & 17 & 2684 & 415 & 80.0 & 11.6 & 6.5 & 1.9 \\
HD$^{34}$S, \cite{hds_ulenikov1995} & 1031.507917 & 16 & 9  & 430  & 130 & --   & --   & --  & --  \\
HD$^{34}$S, наст. работа             & 1031.507611 & 21 & 14 & 1212 & 263 & 85.9 & 9.1  & 2.3 & 2.7 \\
\hline
\end{tabular}
\endgroup

\vspace{0.3em}
\begin{minipage}{0.98\linewidth}
\footnotesize
$N_{tr}$~--- число отнесенных переходов; $N_l$~--- число полученных энергий состояния (010).
Величины $m_i$ показывают доли уровней с остаточными отклонениями:
$m_1$~--- до $10 \times 10^{-5}$~см$^{-1}$,
$m_2$~--- от $10$ до $20 \times 10^{-5}$~см$^{-1}$,
$m_3$~--- от $20$ до $30 \times 10^{-5}$~см$^{-1}$,
$m_4$~--- более $30 \times 10^{-5}$~см$^{-1}$.
\end{minipage}
\end{table}

На основании отнесенных переходов HD$^{32}$S было определено 415 энергий верхнего состояния $(010)$.
Эти значения использовались во взвешенной подгонке параметров эффективного гамильтониана Ватсона в
$A$-редукции. Полученный набор из 33 параметров воспроизводит исходные энергии со
среднеквадратичным отклонением $d_{\textrm{rms}}=1.1\times10^{-4}$ см$^{-1}$. Основные параметры
состояния $(010)$ приведены в табл.~\ref{tab:hds_010_parameters}; распределение остаточных отклонений
показано на рис.~\ref{fig:hds32_residuals}. Видно, что использование расширенного набора линий позволило
устойчиво описать более высокие значения $J$ и $K_a$, чем в предыдущем анализе.

\begin{table}[htbp]
\centering
\caption{Параметры эффективного гамильтониана состояния $(010)$ молекул HD$^{32}$S и HD$^{34}$S,
см$^{-1}$.}
\label{tab:hds_010_parameters}
\scriptsize
\begin{tabular}{l c c}
\hline
Параметр & HD$^{32}$S & HD$^{34}$S \\
\hline
$E$ & 1032.715193(19) & 1031.507615(20) \\
$A$ & 10.0271335(29) & 10.0036871(24) \\
$B$ & 5.0245746(16) & 5.01044009(98) \\
$C$ & 3.19122513(75) & 3.18294464(79) \\
$\Delta_K\times10^3$ & $-0.30303(10)$ & $-0.312897(64)$ \\
$\Delta_{JK}\times10^3$ & 1.042391(54) & 1.043486(51) \\
$\Delta_J\times10^3$ & 0.099853(12) & 0.0988919(45) \\
$\delta_K\times10^3$ & 0.824949(51) & 0.823140(44) \\
$\delta_J\times10^3$ & 0.0350241(59) & 0.0346228(17) \\
$H_K\times10^6$ & 0.8504(15) & 0.83524(79) \\
$H_{KJ}\times10^6$ & $-1.1627(12)$ & $-1.14868(99)$ \\
$H_{JK}\times10^6$ & 0.63952(36) & 0.63493(42) \\
$H_J\times10^6$ & 0.005964(56) & 0.0059599(75) \\
$h_K\times10^6$ & 1.5855(14) & 1.5831(10) \\
$h_{JK}\times10^6$ & 0.30286(31) & 0.30084(18) \\
$h_J\times10^6$ & 0.003135(28) & 0.003135 \\
$L_K\times10^9$ & $-4.162(12)$ & $-4.1204(25)$ \\
$L_{KKJ}\times10^9$ & 5.757(14) & 5.7085(42) \\
$L_{JK}\times10^9$ & $-1.7650(43)$ & $-1.7530(25)$ \\
$L_{JJK}\times10^9$ & $-0.22202(80)$ & $-0.21940(86)$ \\
$L_J\times10^9$ & $-0.001058(88)$ & $-0.001058$ \\
$l_K\times10^9$ & $-2.798(10)$ & $-2.798$ \\
$l_{KJ}\times10^9$ & $-0.6184(60)$ & $-0.6184$ \\
$l_{JK}\times10^9$ & $-0.13201(54)$ & $-0.13201$ \\
$l_J\times10^9$ & $-0.000548(44)$ & $-0.000548$ \\
$P_K\times10^{12}$ & 4.494(61) & 4.494 \\
$P_{KKKJ}\times10^{12}$ & $-6.628(76)$ & $-6.628$ \\
$P_{KKJ}\times10^{12}$ & 2.608(16) & 2.608 \\
$p_K\times10^{12}$ & $-2.191(43)$ & $-2.191$ \\
$p_{KKJ}\times10^{12}$ & $-2.124(19)$ & $-2.124$ \\
$p_{JK}\times10^{12}$ & 0.2155(78) & 0.2155 \\
$Q_K\times10^{15}$ & $-0.831(86)$ & $-0.831$ \\
$Q_{KKKJ}\times10^{15}$ & 0.654(98) & 0.654 \\
$d_{\textrm{rms}}\times10^{-5}$ & 11.0 & 15.9 \\
\hline
\end{tabular}
\par\vspace{1mm}
\footnotesize Значения в скобках соответствуют статистической неопределенности $1\sigma$ в последних знаках.
\end{table}


\begin{figure}[htbp]
    \centering
    \maybeincludegraphics[width=0.95\textwidth]{images/hds_fig4.pdf}
    \caption{Остаточные отклонения <<эксперимент--расчет>> для положений линий, энергий состояния
    $(010)$ и интенсивностей линий полосы $\nu_2$ молекулы HD$^{32}$S.}
    \label{fig:hds32_residuals}
\end{figure}

\subsection{Анализ полосы $\nu_2$ молекулы HD$^{34}$S}

Из-за естественного изотопного состава серы линии HD$^{34}$S значительно слабее линий HD$^{32}$S.
Тем не менее в спектре удалось отнести 1212 переходов полосы $\nu_2$ HD$^{34}$S с максимальными
квантовыми числами $J_{\max}=21$ и $K_{a,\max}=14$. По этим переходам было получено 263 энергии
состояния $(010)$, которые затем использовались для определения параметров эффективного гамильтониана.
Начальные значения параметров брались равными соответствующим значениям HD$^{32}$S. Итоговый набор
параметров воспроизводит исходные энергии HD$^{34}$S со среднеквадратичным отклонением
$d_{\textrm{rms}}=1.6\times10^{-4}$ см$^{-1}$. Остаточные отклонения для положений линий и энергий
состояния $(010)$ HD$^{34}$S представлены на рис.~\ref{fig:hds34_residuals}.

\begin{figure}[htbp]
    \centering
    \maybeincludegraphics[width=0.95\textwidth]{images/hds_fig5.pdf}
    \caption{Остаточные отклонения <<эксперимент--расчет>> для положений линий и энергий состояния
    $(010)$ полосы $\nu_2$ молекулы HD$^{34}$S.}
    \label{fig:hds34_residuals}
\end{figure}

\subsection{Оценка парциального давления HDS в образце}

Как и в случае D$_2$S, температура, полное давление и оптическая длина пути контролировались в ходе
эксперимента, но парциальное давление HDS не могло быть измерено напрямую. Для его определения была
использована теория изотопозамещения~\cite{byk}. В отличие от симметричного замещения
D$_2$S$\leftarrow$H$_2$S, переход HDS$\leftarrow$H$_2$S понижает симметрию молекулы с C$_{2v}$ до C$_s$.
Поэтому при выводе соотношений для параметров дипольного момента необходимо учитывать поворот
молекулярно-фиксированной системы координат. В общем виде связь координатных систем задается
матрицей $K_{\alpha\beta}$:
\begin{equation}
\widetilde{k}_{Z\beta}=\sum_{\alpha} k_{Z\alpha}K_{\alpha\beta},
\label{eq:hds_k_rotation}
\end{equation}
причем элементы этой матрицы зависят от нормальных координат:
\begin{equation}
K_{\alpha\beta}(Q)=K^{e}_{\alpha\beta}+
\sum_{\lambda}K^{\lambda}_{\alpha\beta}Q_{\lambda}+\ldots .
\label{eq:hds_k_expansion}
\end{equation}
С учетом преобразования нормальных координат это приводит к выражению для главных параметров
эффективного дипольного момента изотополога:
\begin{equation}
\widetilde{\mu}^{\nu}_{\gamma}=\sum_{\lambda}\sum_{\alpha}
K^{e}_{\alpha\gamma}
\left(\mu^{\lambda}_{\alpha}-\sum_{\beta}K^{\lambda}_{\alpha\beta}
\widetilde{\mu}^{e}_{\beta}\right)\alpha^{\nu}_{\lambda}.
\label{eq:hds_mu_iso_general}
\end{equation}

На основе параметров H$_2$S и коэффициентов изотопического преобразования были получены теоретические
оценки ${}^{010}\widetilde{\mu}_{x1}=0.01886$~D и
${}^{010}\widetilde{\mu}_{z1}=-0.00855$~D для полосы $\nu_2$ HDS. Эти значения использовались для
расчета интенсивностей 10 изолированных, ненасыщенных и достаточно сильных линий с малыми значениями
$J$ и $K_a$. Затем экспериментальные интенсивности тех же линий были определены из подгонки их
профилей при предположении 100~\% содержания HDS в образце. Парциальное давление рассчитывалось по
формуле
\begin{equation}
P_{\textrm{part}}=\frac{{}^{(\textrm{prob})}S^{N}_{\nu}}{{}^{(\textrm{calc})}S^{N}_{\nu}}
P_{\textrm{samp}},
\label{eq:hds_partial_pressure}
\end{equation}
где $P_{\textrm{samp}}=4.5$ мбар --- полное давление образца. Среднее отношение
${}^{(\textrm{prob})}S^{N}_{\nu}/{}^{(\textrm{calc})}S^{N}_{\nu}$ составило 0.394, откуда было получено
$P_{\textrm{part}}=1.77\pm0.10$ мбар, что соответствует доле HD$^{32}$S около 39.4~\% в исследуемом
образце. Использованные линии и коэффициенты пересчета приведены в табл.~\ref{tab:hds_pressure}.

\begin{table}[htbp]
\centering
\caption{Оценка парциального давления HD$^{32}$S по интенсивностям пробных линий.}
\label{tab:hds_pressure}
\small
\resizebox{\textwidth}{!}{%
\begin{tabular}{ccc ccc c c c c}
\hline
\multicolumn{3}{c}{Верхний уровень} & \multicolumn{3}{c}{Нижний уровень} &
$\nu$, см$^{-1}$ & ${}^{(\textrm{calc})}S^N_\nu$ & ${}^{(\textrm{prob})}S^N_\nu$ &
${}^{(\textrm{calc})}S^N_\nu/{}^{(\textrm{prob})}S^N_\nu$ \\
$J$ & $K_a$ & $K_c$ & $J'$ & $K'_a$ & $K'_c$ & & \multicolumn{2}{c}{10$^{-22}$ см$^{-1}$/(молек. см$^{-2}$)} & \\
\hline
3 & 0 & 3 & 4 & 1 & 4 & 1002.17811 & 6.964 & 2.890 & 2.409 \\
2 & 0 & 2 & 3 & 1 & 3 & 1007.66669 & 5.795 & 2.358 & 2.458 \\
2 & 1 & 2 & 1 & 1 & 1 & 1047.55144 & 1.035 & 0.435 & 2.377 \\
2 & 0 & 2 & 1 & 0 & 1 & 1048.78599 & 1.390 & 0.500 & 2.781 \\
3 & 1 & 3 & 2 & 1 & 2 & 1054.60120 & 1.718 & 0.671 & 2.559 \\
4 & 1 & 4 & 3 & 1 & 3 & 1061.41341 & 2.179 & 0.880 & 2.476 \\
4 & 0 & 4 & 3 & 0 & 3 & 1062.43939 & 2.260 & 0.814 & 2.778 \\
5 & 0 & 5 & 4 & 0 & 4 & 1068.59731 & 2.465 & 0.914 & 2.696 \\
5 & 1 & 5 & 4 & 0 & 4 & 1069.09209 & 9.880 & 4.083 & 2.420 \\
6 & 1 & 6 & 5 & 0 & 5 & 1074.88868 & 10.477 & 4.267 & 2.456 \\
\hline
\multicolumn{9}{r}{Средний коэффициент ${}^{(\textrm{calc})}S^N_\nu/{}^{(\textrm{prob})}S^N_\nu$} & $2.541\pm0.147$ \\
\hline
\end{tabular}%
}
\end{table}


\subsection{Анализ интенсивностей и параметры эффективного дипольного момента}

После определения парциального давления HD$^{32}$S были получены экспериментальные интенсивности
418 линий полосы $\nu_2$ с $J_{\max}=21$ и $K_{a,\max}=6$. Площади линий извлекались из подгонки
профилей Хартмана--Трана~\cite{d2s_tennyson2014}; примеры такой подгонки представлены на
рис.~\ref{fig:hds_lineshape}. Нормированная оптическая толщина записывалась в виде
\begin{equation}
\tau(\nu)^{\textrm{exp}}=\frac{1}{\lg e}\lg\left(\frac{I_0(\nu)}{I(\nu)}\right).
\label{eq:hds_tau}
\end{equation}

\begin{figure}[htbp]
    \centering
    \maybeincludegraphics[width=0.86\textwidth]{images/hds_fig6.pdf}
    \caption{Примеры подгонки профилей отдельных линий полосы $\nu_2$ молекулы HD$^{32}$S с
    использованием профиля Хартмана--Трана.}
    \label{fig:hds_lineshape}
\end{figure}

Экспериментальные интенсивности были использованы для подгонки параметров эффективного оператора
дипольного момента. Для HD$^{32}$S был получен набор параметров, приведенный в
табл.~\ref{tab:hds_dipole_parameters}. Он воспроизводит 418 исходных интенсивностей со
среднеквадратичным отклонением $d_{\textrm{rms}}=4.7$~\%. Отношение главных параметров
${}^{010}\widetilde{\mu}_{x1}$ и ${}^{010}\widetilde{\mu}_{z1}$, найденное из подгонки, хорошо
согласуется с теоретической оценкой: экспериментальное значение составляет 2.183, а расчетное --- 2.20.

\begin{table}[htbp]
\centering
\caption{Параметры эффективного оператора дипольного момента полосы $\nu_2$ молекул HDS, D.}
\label{tab:hds_dipole_parameters}

\begingroup
\small
\setlength{\tabcolsep}{6pt}
\renewcommand{\arraystretch}{1.05}
\begin{tabular*}{\textwidth}{@{}l@{\extracolsep{\fill}}lcc@{}}
\hline
Оператор & Параметр & HD$^{32}$S & HD$^{34}$S \\
\hline
$k_{zx}$ &
${}^{010}\mu_{x1}\times10^{-1}$ &
0.19615(45) &
0.19616(96) \\

$\{ik_{zy},J_z\}$ &
${}^{010}\mu_{x4}\times10^{-3}$ &
0.6673(99) &
0.6673 \\

$\{k_{zz},iJ_y\}$ &
${}^{010}\mu_{x5}\times10^{-4}$ &
0.407(42) &
0.407 \\

$k_{zz}$ &
${}^{010}\mu_{z1}\times10^{-2}$ &
$-0.8987(88)$ &
$-0.853$ \\

$\frac{1}{2}\left[\{k_{zx},iJ_y\}-\{ik_{zy},J_x\}\right]$ &
${}^{010}\mu_{z4}\times10^{-3}$ &
0.453(20) &
0.453 \\

$\frac{1}{2}\left[\{k_{zx},iJ_y\}+\{ik_{zy},J_x\}\right]$ &
${}^{010}\mu_{z6}\times10^{-3}$ &
$-0.598(34)$ &
$-0.598$ \\
\hline
\end{tabular*}
\endgroup

\vspace{0.3em}
\begin{minipage}{0.98\linewidth}
\footnotesize
Здесь $\{A,B\}=AB+BA$, $J^2_{xy}=J_x^2-J_y^2$.
Значения в скобках соответствуют стандартным ошибкам $1\sigma$.
\end{minipage}
\end{table}

Для HD$^{34}$S был выполнен аналогичный, но более ограниченный анализ интенсивностей. Из-за меньшей
естественной распространенности этого изотополога надежно выделяются только наиболее сильные линии;
для них сравнение рассчитанных и экспериментальных интенсивностей дает среднеквадратичное отклонение
около 4.0~\%. В итоговый список были включены положения и рассчитанные интенсивности 2684 переходов
HD$^{32}$S и 1212 переходов HD$^{34}$S полосы $\nu_2$. Полученный набор параметров обеспечивает
согласованное описание как положений линий, так и их абсолютных интенсивностей для двух
монодейтерированных изотопологов сероводорода.

\section*{Выводы по главе 3}
\addcontentsline{toc}{section}{Выводы по главе 3}

В результате проведенного исследования получены следующие основные результаты:
\begin{enumerate}
    \item В спектре фундаментальной полосы $\nu_2$ молекулы D$_2{}^{32}$S отнесено 1742 перехода с
    $J_{\max}=30$ и $K_{a,\max}=21$. По ним определено 535 энергий состояния $(010)$ и получен набор из
    34 параметров эффективного гамильтониана, воспроизводящий исходные энергии со среднеквадратичным
    отклонением $1.31\times10^{-4}$ см$^{-1}$.

    \item На основе теории изотопозамещения определено парциальное давление D$_2{}^{32}$S в образце и
    выполнен анализ 280 экспериментальных линий, соответствующих 324 переходам. Набор из восьми
    параметров эффективного дипольного момента воспроизводит измеренные интенсивности со
    среднеквадратичным отклонением 2.5~\%.

    \item Для полосы $\nu_2$ молекулы HD$^{32}$S отнесено 2684 перехода и определено 415 энергий
    состояния $(010)$; полученный набор из 33 параметров воспроизводит их со среднеквадратичным
    отклонением $1.1\times10^{-4}$ см$^{-1}$. Для HD$^{34}$S отнесено 1212 переходов, определено 263
    энергетических значения, а среднеквадратичное отклонение их описания составило
    $1.6\times10^{-4}$ см$^{-1}$.

    \item Для HD$^{32}$S определено парциальное давление в газовой смеси и проанализированы
    интенсивности 418 линий полосы $\nu_2$. Полученные параметры эффективного дипольного момента
    воспроизводят исходные интенсивности со среднеквадратичным отклонением 4.7~\%; для наиболее
    сильных линий HD$^{34}$S соответствующее отклонение составляет около 4.0~\%.
\end{enumerate}

Таким образом, для трех дейтерированных изотопологов сероводорода сформированы согласованные наборы
параметров, обеспечивающие описание энергетической структуры и абсолютных интенсивностей линий
фундаментальной полосы $\nu_2$ с точностью, близкой к экспериментальной. Полученные списки переходов
могут использоваться при моделировании высокоразрешенных спектров, уточнении поверхностей
потенциальной энергии и дипольного момента, а также при подготовке спектроскопических данных для
прикладных задач атмосферной оптики и астрофизики.

